\documentclass[11pt]{article}
\usepackage[margin=1in]{geometry}
\usepackage{amsmath,amssymb,amsthm}
\usepackage{parskip}
\usepackage{booktabs}
\usepackage{array}

\newtheorem{theorem}{Theorem}
\newtheorem{lemma}[theorem]{Lemma}
\newtheorem{corollary}[theorem]{Corollary}
\newtheorem{proposition}[theorem]{Proposition}
\theoremstyle{definition}
\newtheorem{definition}[theorem]{Definition}
\theoremstyle{remark}
\newtheorem{remark}[theorem]{Remark}

\newcommand{\Z}{\mathbb{Z}}
\newcommand{\R}{\mathbb{R}}
\newcommand{\C}{\mathbb{C}}
\newcommand{\tauc}{\tau}

\title{Disproof of conjecture \texttt{00000000462}}
\author{earthking11}
\date{2026-10-04}

\begin{document}
\maketitle

\section*{Verdict: FALSE}

We disprove conjecture \texttt{00000000462} as stated. Let $\tauc(n)$ denote the
number of spanning trees of the circulant graph $C_n(1,2)$, the \emph{square of
the $n$-cycle}: the graph on $\Z/n$ joining $i$ to $i\pm 1$ and to $i\pm 2$. We
compute $\tauc$ exactly, and show:

\begin{enumerate}
\item $\tauc$ satisfies the six-term linear recurrence
\[
\tauc(n+6) \;=\; 4\,\tauc(n+5) - 10\,\tauc(n+3) + 4\,\tauc(n+1) - \tauc(n),
\]
whose characteristic polynomial is
\[
p(x) \;=\; x^{6} - 4x^{5} + 10x^{3} - 4x + 1 \;=\; \bigl((x+1)(x^{2}-3x+1)\bigr)^{2}.
\]
Its largest real root is
\[
\frac{3+\sqrt5}{2} \;=\; \varphi^{2} \;=\; 2.6180339887\ldots,
\qquad \varphi = \frac{1+\sqrt5}{2}.
\]
\item The conjectured value
\[
\alpha^{2} \;=\; (2+\sqrt3)^{2} \;=\; 7 + 4\sqrt3 \;=\; 13.9282032302\ldots
\]
is \emph{not} a root of $p$: an exact computation in $\Z[\sqrt3]$ gives
\[
p\bigl(7+4\sqrt3\bigr) \;=\; 2615536 + 1510080\sqrt3 \;\neq\; 0 .
\]
\end{enumerate}

Consequently $\alpha^{2}$ cannot be a characteristic root of the recurrence
satisfied by $\tauc$, and it is also not the exponential growth rate, which is
$\varphi^{2}$. The conjecture fails under either reading of its wording; see
Section~\ref{sec:readings}.

\section{The conjecture}

Quoted verbatim from \texttt{conjectures/00000000462.md}:

\begin{quote}
\textbf{English.} Definition: $\tauc(\text{circulant } C_n(1,2))$. Conjecture:
$\tauc(C_n(1,2))$ satisfies a linear recurrence whose largest real
characteristic root tends to $\alpha^{2}$, where $\alpha = 2+\sqrt3$ (directly
verifiable against the matrix-tree theorem, then stateable in closed form).
\end{quote}

\noindent The Chinese statement filed with the conjecture is reproduced verbatim
in \texttt{README.md} together with this English text; it is not typeset here so
that \texttt{main.tex} has no CJK font requirement and compiles with a plain
\texttt{tectonic main.tex}. The two formulations are translations of one
another, and both are refuted below.

\begin{definition}[The graph]
For $n \ge 5$ let $C_n(1,2)$ be the circulant graph on $\Z/n\Z$ in which $i$ and
$j$ are adjacent whenever $j-i \equiv \pm 1$ or $j-i \equiv \pm 2 \pmod n$. Let
$\tauc(n) = \tauc(C_n(1,2))$ be its number of spanning trees.
\end{definition}

\begin{remark}
For $n = 3$ and $n = 4$ the connection set $\{\pm1,\pm2\}$ collapses
($\pm2 = \mp1$ for $n=3$, and $\pm2 = \{2\}$ with $\pm1 = \{1,3\}$ for $n=4$),
and the graphs are the complete graphs $K_3$ and $K_4$ with $\tauc(3)=3$ and
$\tauc(4)=16$. The closed form below is the one for the general case $n \ge 5$,
which is the case the conjecture is about. Every statement we prove is checked
numerically for $5 \le n \le 60$.
\end{remark}

\section{The exact spanning-tree count}

\begin{theorem}[Closed form]
\label{thm:closed}
For every $n \ge 5$,
\[
\tauc(n) \;=\; \frac{n\bigl(L(2n) + 2(-1)^{n-1}\bigr)}{5},
\]
where $L$ is the Lucas sequence $L(0)=2$, $L(1)=1$,
$L(k+2)=L(k+1)+L(k)$. In particular $\tauc(n)$ is a positive integer and
$5 \mid L(2n) + 2(-1)^{n-1}$.
\end{theorem}

\begin{proof}
By the matrix-tree theorem (Kirchhoff), for a connected graph on $n$ vertices
with Laplacian eigenvalues $\lambda_0 = 0, \lambda_1, \dots, \lambda_{n-1}$,
\[
\tauc \;=\; \frac1n \prod_{k=1}^{n-1} \lambda_k .
\]
For a circulant graph with connection set $S$, the Laplacian eigenvalues are
$\lambda_k = \sum_{s \in S} \bigl(1 - \omega^{ks}\bigr)$ with
$\omega = e^{2\pi i/n}$. Here $S = \{\pm 1, \pm 2\}$, so with $z = \omega^{k}$,
\begin{equation}
\label{eq:lambda}
\lambda_k \;=\; 4 - z - z^{-1} - z^{2} - z^{-2}
\;=\; -\frac{(z-1)^{2}\,(z^{2}+3z+1)}{z^{2}},
\end{equation}
the factorisation following from
$z^{4}+z^{3}-4z^{2}+z+1 = (z-1)^{2}(z^{2}+3z+1)$.

Set $z_k = \omega^{k}$ for $k = 1, \dots, n-1$; these are the nontrivial $n$-th
roots of unity, and $\prod_{k=1}^{n-1} z_k^{2} = \omega^{2 \cdot n(n-1)/2} = 1$.
Also
\[
\prod_{k=1}^{n-1}(z_k - 1) \;=\; (-1)^{n-1}\prod_{k=1}^{n-1}(1-\omega^{k}) \;=\; (-1)^{n-1} n ,
\]
because $\prod_{k=1}^{n-1}(x - \omega^{k}) = (x^{n}-1)/(x-1)$ evaluated at
$x = 1$, whence
\begin{equation}
\label{eq:prodone}
\prod_{k=1}^{n-1}(z_k-1)^{2} \;=\; n^{2}.
\end{equation}
Let $r_1, r_2$ be the roots of $z^{2}+3z+1$, so $r_1 + r_2 = -3$ and
$r_1 r_2 = 1$. Since
$\prod_{k=1}^{n-1}(x - z_k) = (x^{n}-1)/(x-1)$ for $x \notin \{z_k\}$,
\[
\prod_{k=1}^{n-1}(z_k - r) \;=\; (-1)^{n-1}\,\frac{r^{n}-1}{r-1},
\]
hence
\begin{equation}
\label{eq:prodq}
\prod_{k=1}^{n-1}(z_k^{2}+3z_k+1)
= \prod_{k=1}^{n-1}(z_k - r_1)(z_k - r_2)
= \frac{(r_1^{n}-1)(r_2^{n}-1)}{(r_1-1)(r_2-1)}
= \frac{(r_1^{n}-1)(r_2^{n}-1)}{5},
\end{equation}
because $(r_1-1)(r_2-1) = r_1r_2 - (r_1+r_2) + 1 = 1+3+1 = 5$. Combining
\eqref{eq:lambda}, \eqref{eq:prodone} and \eqref{eq:prodq},
\[
\tauc(n) \;=\; \frac1n \prod_{k=1}^{n-1}\lambda_k
\;=\; \frac1n \cdot (-1)^{n-1} \cdot n^{2} \cdot \frac{(r_1^{n}-1)(r_2^{n}-1)}{5}
\;=\; \frac{(-1)^{n-1}\,n\,(r_1^{n}-1)(r_2^{n}-1)}{5}.
\]
Now $(r_1^{n}-1)(r_2^{n}-1) = (r_1r_2)^{n} - r_1^{n} - r_2^{n} + 1
= 2 - (r_1^{n}+r_2^{n})$. Writing
$s_1 = \tfrac{3-\sqrt5}{2}$ and $s_2 = \tfrac{3+\sqrt5}{2}$ we have
$r_1 = -s_1$, $r_2 = -s_2$, so
$r_1^{n}+r_2^{n} = (-1)^{n}(s_1^{n}+s_2^{n})$. Since $s_1, s_2$ are the roots of
$x^{2}-3x+1$ with $s_1^{0}+s_2^{0} = 2$ and $s_1+s_2 = 3$, the sequence
$s_1^{n}+s_2^{n}$ is $L(2n)$: indeed $L(2n)$ satisfies the same recurrence
$x^{2}-3x+1$ with initial values $L(0)=2$, $L(2)=3$. Therefore
\[
(r_1^{n}-1)(r_2^{n}-1) \;=\; 2 - (-1)^{n} L(2n),
\]
and
\[
\tauc(n) \;=\; \frac{(-1)^{n-1} n \bigl(2 - (-1)^{n}L(2n)\bigr)}{5}
\;=\; \frac{n\bigl(L(2n) + 2(-1)^{n-1}\bigr)}{5},
\]
since $-(-1)^{n-1}(-1)^{n} = +1$. This is the asserted form.
\end{proof}

The integrality statement in Theorem~\ref{thm:closed} is an instance of a small
congruence: $L(2n) \equiv 2 \pmod 5$ for even $n$ and $L(2n) \equiv 3 \pmod 5$
for odd $n$, so $L(2n) + 2(-1)^{n-1} \equiv 2 - 2 \equiv 0$ (even $n$) and
$\equiv 3 + 2 \equiv 0$ (odd $n$) modulo $5$.

\section{The linear recurrence and its characteristic polynomial}

Write $u(n) = L(2n) + 2(-1)^{n-1}$, so that $\tauc(n) = n\,u(n)/5$ by
Theorem~\ref{thm:closed}.

\begin{lemma}[The recurrence for $u$]
\label{lem:urec}
For every $n \ge 0$,
\[
u(n+3) \;=\; 2\,u(n+2) + 2\,u(n+1) - u(n).
\]
Equivalently, $u$ satisfies the linear recurrence with characteristic
polynomial
\[
q(x) \;=\; x^{3} - 2x^{2} - 2x + 1 \;=\; (x+1)\,(x^{2}-3x+1).
\]
\end{lemma}

\begin{proof}
The even-index Lucas numbers satisfy
$L(2n+4) = 3L(2n+2) - L(2n)$; this is the standard Lucas identity
$L(m+4) + L(m) = 3L(m+2)$ read at $m = 2n$, and it is the recurrence with
characteristic polynomial $x^{2}-3x+1$. The sequence $(-1)^{n-1}$ satisfies
$x+1$. Summing the two recurrences in operator form gives
\[
(E^{2}-3E+1)(E+1)\,u \;=\; 0,
\]
i.e.\ $(E^{3}-2E^{2}-2E+1)u = 0$, which is the stated three-term recurrence.
The factorisation $x^{3}-2x^{2}-2x+1 = (x+1)(x^{2}-3x+1)$ is verified by
expansion: $(x+1)(x^{2}-3x+1) = x^{3}-3x^{2}+x + x^{2}-3x+1 = x^{3}-2x^{2}-2x+1$.
\end{proof}

\begin{lemma}[The recurrence for $\tauc$]
\label{lem:taurec}
For every $n \ge 0$ the sequence $\tauc$ satisfies
\[
\tauc(n+6) \;=\; 4\,\tauc(n+5) - 10\,\tauc(n+3) + 4\,\tauc(n+1) - \tauc(n),
\]
so its characteristic polynomial divides
\[
p(x) \;=\; q(x)^{2} \;=\; x^{6} - 4x^{5} + 10x^{3} - 4x + 1
\;=\; \bigl((x+1)(x^{2}-3x+1)\bigr)^{2}.
\]
\end{lemma}

\begin{proof}
Multiplying a sequence satisfying $q(E)w = 0$ by $n$ doubles the multiplicity of
every root: if $q(E)w = 0$ then $q(E)^{2}(n\,w) = 0$, because
$q(E)\bigl(n\,w(n)\bigr)$ is a fixed finite linear combination of the shifts of
$w$, which $q(E)$ annihilates. Applying this to $w = u$ and using
$\tauc(n) = n\,u(n)/5$ gives $p(E)\tauc = 0$ with $p = q^{2}$.

For $q(x) = x^{3}-2x^{2}-2x+1$ we compute
\[
q(x)^{2} = (x^{3}-2x^{2}-2x+1)^{2} = x^{6} - 4x^{5} + 10x^{3} - 4x + 1 .
\]
Read as an operator identity $E^{6} - 4E^{5} + 10E^{3} - 4E + 1 = 0$, this is
exactly the displayed recurrence. (The computation was also verified
symbolically and numerically; see \texttt{reproduce.py}.)
\end{proof}

\begin{theorem}[Roots]
\label{thm:roots}
The roots of $p$ are $-1$ (twice) and
\[
s_2 = \frac{3+\sqrt5}{2} = \varphi^{2} \approx 2.6180339887,
\qquad
s_1 = \frac{3-\sqrt5}{2} \approx 0.3819660113,
\]
each twice (the roots of $x^{2}-3x+1$). The largest real root is $s_2 = \varphi^{2}$
and it is simple as a root of $p$ in the sense that its multiplicity as a root
of the minimal annihilator of $\tauc$ is at least the multiplicity forced by the
factor $n$; the exponential growth rate of $\tauc$ is $\varphi^{2}$.
\end{theorem}

\begin{proof}
Immediate from $p = ((x+1)(x^{2}-3x+1))^{2}$ and the quadratic formula. The
growth statement follows from Theorem~\ref{thm:closed}: $\tauc(n) = n\,u(n)/5$
and $u(n) = s_1^{n} + s_2^{n} + 2(-1)^{n-1}$ with $|s_1| < 1 < s_2$, so
$\tauc(n) = \bigl(s_2^{n}(1+o(1))\bigr)\cdot n/5$ and
$\lim_{n\to\infty} \tauc(n)^{1/n} = s_2 = \varphi^{2}$.
\end{proof}

\section{The conjectured root is not a characteristic root}
\label{sec:root}

\begin{theorem}
\label{thm:notroot}
$\alpha^{2} = (2+\sqrt3)^{2} = 7+4\sqrt3$ is not a root of
$p(x) = x^{6}-4x^{5}+10x^{3}-4x+1$.
\end{theorem}

\begin{proof}
Because $p(x) = (x+1)^{2}(x^{2}-3x+1)^{2}$, a number is a root of $p$ iff it
equals $-1$ or is a root of $x^{2}-3x+1$. Work in the ring
$\Z[\sqrt3] \cong \Z[t]/(t^{2}-3)$, embedded in $\R$ by $t \mapsto \sqrt3$.

With $\alpha = 2+\sqrt3$ we get $\alpha^{2} = 7+4\sqrt3$ and
$\alpha^{4} = (\alpha^{2})^{2} = 97 + 56\sqrt3$ in $\Z[\sqrt3]$. Hence
\[
\alpha^{4} - 3\alpha^{2} + 1
= (97+56\sqrt3) - 3(7+4\sqrt3) + 1
= 77 + 44\sqrt3 .
\]
Since $77 > 0$ and $44\sqrt3 > 0$, this number is nonzero. Also
$\alpha^{2} = 7+4\sqrt3 > 0 > -1$. Therefore
\[
p(\alpha^{2}) \;=\; (\alpha^{2}+1)^{2}\bigl(\alpha^{4}-3\alpha^{2}+1\bigr)^{2}
\;=\; (8+4\sqrt3)^{2}\,(77+44\sqrt3)^{2} \;>\; 0 ,
\]
and expanding, $p(\alpha^{2}) = 2615536 + 1510080\sqrt3 \neq 0$.
\end{proof}

\begin{corollary}[Refutation]
\label{cor:refute}
The sequence $\tauc(C_n(1,2))$ satisfies a linear recurrence with
characteristic polynomial $p$, and the conjectured value
$\alpha^{2} = 7+4\sqrt3$ is not a root of $p$. Since the minimal annihilating
polynomial of a linear recurrence divides every annihilating polynomial, every
characteristic root of $\tauc$ is a root of $p$; hence $\alpha^{2}$ is not a
characteristic root of any linear recurrence satisfied by $\tauc$. The
conjecture is false as stated.
\end{corollary}

\begin{proof}
Lemma~\ref{lem:taurec} exhibits the annihilating polynomial $p$ of $\tauc$;
Theorem~\ref{thm:notroot} shows $\alpha^{2}$ is not a root of $p$. If $q$ is the
minimal annihilating polynomial of $\tauc$ then $q \mid p$, so every root of $q$
is a root of $p$; since $p(\alpha^{2}) \neq 0$ we get $q(\alpha^{2}) \neq 0$.
\end{proof}

\section{The two readings of the conjecture}
\label{sec:readings}

The filed wording says that $\tauc$ ``satisfies a linear recurrence whose
largest real characteristic root \emph{tends to} $\alpha^{2}$''. The word
``tends'' is not standard for characteristic roots of a constant-coefficient
recurrence; we dispose of both natural readings.

\begin{enumerate}
\item \textbf{``Characteristic root'' $=$ root of the characteristic
polynomial.} For a linear recurrence $a_{n+k} = c_{k-1}a_{n+k-1} + \dots +
c_0a_n$ the characteristic roots are the roots of
$x^{k} - c_{k-1}x^{k-1} - \dots - c_0$, i.e.\ the roots of any annihilating
polynomial. By Corollary~\ref{cor:refute} every such root of $\tauc$ is a root
of $p$, and $\alpha^{2}$ is not one. \emph{The conjecture is false.}

\item \textbf{``Tends to'' $=$ the exponential growth rate is $\alpha^{2}$.}
If the claim is that $\lim_{n\to\infty}\tauc(n)^{1/n} = \alpha^{2}$, then the
limit is in fact $\varphi^{2} = 2.6180339887\ldots$ by Theorem~\ref{thm:roots},
not $\alpha^{2} = 13.9282032302\ldots$. Indeed $\tauc(n) \sim n\,\varphi^{2n}/5$:
for $n = 60$ the ratio
$\frac{\tauc(60)/60}{\tauc(59)/59}$ equals $2.6180339887$ to ten decimal places.
\emph{The conjecture is false.}
\end{enumerate}

\begin{remark}[Why the numbers differ]
$\varphi^{2} = \frac{3+\sqrt5}{2}$ comes from the Lucas recurrence
$x^{2}-3x+1$, which is the recurrence genuinely satisfied by the even-index
Lucas numbers underlying $u(n) = L(2n)+2(-1)^{n-1}$. The filed value
$\alpha^{2} = 7+4\sqrt3$ comes from the order-$2$ algebraic number
$2+\sqrt3$; it is not related to $C_n(1,2)$ in any way we can detect, and
$p(\alpha^{2}) \neq 0$ rules it out as a root of the annihilator. Note also that
$\alpha = 2+\sqrt3$ is \emph{not} of the form $\varphi^{k}$ for integer $k$
($\varphi^{3} = 2+\sqrt5$), so it cannot arise from the Lucas family either.
\end{remark}

\section{Numerical verification}

The following table collects the first values. The ``matrix tree'' column is the
exact integer determinant of the reduced Laplacian (fraction-free Bareiss
elimination); the ``closed form'' column is
$n\bigl(L(2n)+2(-1)^{n-1}\bigr)/5$.

\begin{center}
\begin{tabular}{r r r r}
\toprule
$n$ & matrix tree $\tauc(n)$ & closed form & recurrence residual \\
\midrule
$5$  & $125$      & $125$      & $0$ \\
$6$  & $384$      & $384$      & $0$ \\
$7$  & $1183$     & $1183$     & $0$ \\
$8$  & $3528$     & $3528$     & $0$ \\
$9$  & $10404$    & $10404$    & $0$ \\
$10$ & $30250$    & $30250$    & $0$ \\
$11$ & $87131$    & $87131$    & $0$ \\
$12$ & $248832$   & $248832$   & $0$ \\
\bottomrule
\end{tabular}
\end{center}

\noindent The residual column is
$\tauc(n) - \bigl(4\tauc(n-1) - 10\tauc(n-3) + 4\tauc(n-5) - \tauc(n-6)\bigr)$.

The two candidate roots compare as follows.

\begin{center}
\begin{tabular}{l r l}
\toprule
quantity & value & status \\
\midrule
$\varphi^{2} = (3+\sqrt5)/2$ & $2.6180339887\ldots$ & root of $p$ (the actual growth rate) \\
$\alpha^{2} = 7+4\sqrt3$    & $13.9282032302\ldots$ & \textbf{not} a root: $p(\alpha^{2}) = 2615536 + 1510080\sqrt3$ \\
\bottomrule
\end{tabular}
\end{center}

\section{Formalisation}

The refutation is formalised in the Lean~4 project under \texttt{lean4/}
(core Lean only --- no Mathlib, no \texttt{sorry}, no \texttt{native\_decide}).
The project defines the Lucas numbers, the sequence
$u(n) = L(2n)+2(-1)^{n-1}$, proves $u$'s three-term recurrence, derives the
six-term recurrence for $c(n) = n\,u(n) = 5\tauc(n)$, proves $c$ and $\tauc$
satisfy it, proves $5 \mid u(n)$, verifies the initial values
$\tauc(5),\dots,\tauc(12)$, verifies the factorisation of $p$ by exact
coefficient arithmetic, models $\Z[\sqrt3]$ as a pair ring, and checks
$p(\alpha^{2}) \neq 0$. The main formal statements are

\begin{quote}\ttfamily
theorem u\_rec (n) : u (n+3) = 2*u(n+2) + 2*u(n+1) - u n\\
theorem tau\_rec (n) : tau (n+6) = 4*tau(n+5) - 10*tau(n+3) + 4*tau(n+1) - tau n\\
theorem char\_square : pmul [1,-2,-2,1] [1,-2,-2,1] = [1,-4,0,10,0,-4,1]\\
theorem pEval\_alpha\_sq : pEval (mul alpha alpha) = (2615536, 1510080)\\
theorem conjecture\_00000000462\_false :\\
\hspace*{2em}(forall n, tau (n+6) = 4*tau(n+5) - 10*tau(n+3) + 4*tau(n+1) - tau n)\\
\hspace*{2em}and (forall x, x = mul alpha alpha -> pEval x != zero)\\
\hspace*{2em}and mul alpha alpha != scale (-1) one
\end{quote}

\noindent The axiom audit (\texttt{lake env lean Check.lean}) reports only
\texttt{propext} and \texttt{Quot.sound} for every theorem, and no
\texttt{sorryAx}. See \texttt{lean4/README.md} for the full statement table and
scope note.

\section{Reproducibility}

\begin{quote}\ttfamily
\$ python3 reproduce.py\\
... (matrix-tree determinant for 5 <= n <= 60, the closed form,\\
... the recurrence, the factorisation, and the exact check in Z[sqrt 3])\\
PASS: every check verified; conjecture 00000000462 is FALSE.
\end{quote}

\texttt{reproduce.py} uses only the Python~3 standard library and exact integer
arithmetic throughout (fraction-free Bareiss elimination for the determinant,
integer coefficient lists for polynomials, pairs for $\Z[\sqrt3]$ and
$\Z[\sqrt5]$); it runs in under a second, prints \texttt{PASS} and exits $0$
exactly when every check holds. The Lean project is built and audited by

\begin{quote}\ttfamily
cd lean4\\
lake build\\
lake env lean Check.lean
\end{quote}

\noindent and the document is compiled by

\begin{quote}\ttfamily
tectonic main.tex\\
(equivalently: tectonic --outdir build main.tex)
\end{quote}

\noindent producing \texttt{build/main.pdf}.

\section{Status against the submission rules}

\begin{itemize}
\item \textbf{LaTeX source} --- present (\texttt{main.tex}), a standalone
\texttt{article} using only \texttt{geometry}, \texttt{amsmath},
\texttt{amssymb}, \texttt{amsthm}, \texttt{parskip}, \texttt{booktabs}, and
\texttt{array}; it compiles with \texttt{tectonic main.tex}.
\item \textbf{PDF document} --- present at \texttt{build/main.pdf}.
\item \textbf{Lean 4 project} --- present under \texttt{lean4/}, with
\texttt{lean-toolchain} (pinned to \texttt{leanprover/lean4:v4.33.1}),
\texttt{lakefile.toml} (library \texttt{Main}, project \texttt{tlmc462}),
\texttt{Main.lean} (core Lean only, no \texttt{sorry}, no
\texttt{native\_decide}), \texttt{Check.lean} (\texttt{\#print axioms} audit),
and \texttt{README.md}. The audit reports only \texttt{[propext, Quot.sound]}
and no \texttt{sorryAx}.
\item \textbf{Auxiliary code} --- \texttt{reproduce.py}, dependency-free, exits
$0$ and prints \texttt{PASS}.
\item \textbf{No other files touched} --- the pull request adds only the folder
\texttt{solutions/00000000462/earthking11\_submission\_20261004212700/} and
modifies nothing else.
\end{itemize}

\end{document}
